% ============================================================================
%  Section IV — Sensor and Signal Characterization  (IEEE TIM, IEEEtran)
%  Passive voice / third person throughout. Requires: graphicx.
%  Figures: fig_psd_usable_band.png, fig_channel_mi.png (in results/).
%  Numbers from results/analysis_B.{md,json}.
% ============================================================================

\section{Sensor and Signal Characterization}
\label{sec:charact}

Before any classifier is trained, it is instructive to treat the smartphone
accelerometer as a measurement instrument and to establish what fault-related
information its channels can, in principle, convey. The metrological
characterisation of MEMS accelerometers---their resolution, noise, calibration,
and traceability---has itself become an active subject in the measurement
community~\cite{prato2021mems, schiavi2023mems, demilia2021mems}, and the same
instrument-oriented perspective is adopted here for a consumer-grade device used
in condition monitoring. This section characterises the effective resolution and
noise floor of the device, delimits the usable frequency band imposed by the
$100$~Hz sampling rate, and quantifies the relative diagnostic value of the raw
and gravity-compensated channels. All quantities were computed directly from the
recordings, without any model training.

\subsection{Effective Resolution and Noise Floor}
The effective amplitude resolution was estimated as the smallest positive
difference between consecutive distinct sample values in a healthy, unloaded
recording. A quantisation step of approximately $0.015$~mg was obtained for the
raw acceleration channels and $0.001$~mg for the gravity-compensated channels,
the finer step of the latter reflecting the additional numerical processing of
the on-board sensor fusion. The broadband noise floor was assessed from the
power spectral density (PSD) of a healthy recording; the median PSD above
$30$~Hz was $8.2\times10^{-8}$~g$^2$/Hz ($-70.8$~dB). Because the fault-induced
vibration reaches amplitudes on the order of $0.1$--$0.2$~g
(Fig.~\ref{fig:signals}), which exceeds this floor by several orders of
magnitude, the resolution and noise characteristics of the consumer sensor are
not the limiting factor for the coarse fault-discrimination task considered here.

\subsection{Usable Frequency Band}
At a sampling rate of $100$~Hz the analysable spectrum is bounded by the Nyquist
frequency of $50$~Hz~\cite{nyquist1924certain, shannon1948mathematical}.
Fig.~\ref{fig:psd} shows the Welch PSD of the gravity-compensated $g_{userx}$
channel for the six health states under identical operating conditions. Two
observations follow. First, discriminative structure is present within the
$0$--$50$~Hz band: the healthy machine exhibits a spectrum roughly an order of
magnitude lower in power than the fault conditions, and distinct spectral peaks
appear for the rotor and electrical faults. The dominant peaks near $24$~Hz
coincide with the shaft rotational frequency at $50$~Hz supply
(Table~\ref{tab:speeds}, e.g.\ $1425$--$1486$~r/min $\approx 23.8$--$24.8$~Hz),
confirming that the low-frequency band captures the rotational fundamental and
its low-order harmonics. Second, the characteristic bearing defect frequencies
(ball-pass and ball-spin frequencies and their harmonics), which for this bearing
and speed range lie well above $50$~Hz, cannot be represented. The dataset is
therefore suited to time-domain and low-frequency analysis rather than the
high-frequency envelope analysis conventionally used for incipient bearing
diagnosis, and a corresponding difficulty in separating the fine
lubrication-severity states is anticipated.

\subsection{Raw versus Gravity-Compensated Channels}
The dataset provides both the raw acceleration and the gravity-compensated
(\textit{userAcceleration}) triads, and their relative diagnostic value is not
self-evident. For each axis, the Pearson correlation between the mean-removed raw
signal and the corresponding gravity-compensated signal was computed over every
recording; the mean correlations were $0.58$, $0.46$, and $0.46$ for the $x$,
$y$, and $z$ axes, respectively. This moderate correlation indicates that the
on-device gravity compensation performs more than a simple removal of the static
component: its complementary filter also attenuates low-frequency content that is
shared by the two representations. To quantify diagnostic relevance, the mutual
information between each channel's time-domain feature set and the fault label was
estimated (Fig.~\ref{fig:mi}). The raw channels were consistently more informative
(mean mutual information $0.65$--$0.75$) than the gravity-compensated channels
($0.46$--$0.57$), indicating that the low-frequency content discarded by the
compensation filter carries useful fault information. Two conclusions of practical
importance for smartphone-based acquisition follow: the raw channels are
preferable for condition monitoring, and the two triads are complementary rather
than redundant. This informativeness ordering also motivates the channel-selection
study reported in Section~\ref{sec:results}.

% ---------------------------------------------------------------------------
\begin{figure}[t]
\centering
\includegraphics[width=\columnwidth]{fig_psd_usable_band.png}
\caption{Welch power spectral density of the $g_{userx}$ channel for the six
health states (loaded, $50$~Hz supply). The analysable band is bounded by the
$50$~Hz Nyquist frequency; the healthy state is markedly lower in power, and
rotor/electrical faults produce distinct low-frequency peaks.}
\label{fig:psd}
\end{figure}

% ---------------------------------------------------------------------------
\begin{figure}[t]
\centering
\includegraphics[width=\columnwidth]{fig_channel_mi.png}
\caption{Mean mutual information between each channel's time-domain feature set
and the fault label. The raw channels ($g_x$, $g_y$, $g_z$) are more informative
than the gravity-compensated channels ($g_{userx}$, $g_{usery}$, $g_{userz}$).}
\label{fig:mi}
\end{figure}
